% !TEX program = lualatex
% Generated from Markdown by src/core/latex/export.ts. Edit the Markdown source, not this file.
\documentclass[11pt,a4paper]{article}
\usepackage{iftex}
\ifPDFTeX
  \usepackage[T1]{fontenc}
  \usepackage[utf8]{inputenc}
  \usepackage{lmodern}
\else
  \usepackage{fontspec}
\fi
\usepackage[brazilian,main=english]{babel}
\usepackage{csquotes}
\usepackage{mathtools}
\usepackage{amssymb}
\usepackage{amsthm}
\usepackage{graphicx}
\usepackage{booktabs}
\usepackage{enumitem}
\usepackage[normalem]{ulem}
\usepackage{fancyvrb}
\usepackage{newfloat}
\usepackage{caption}
% ABNT NBR 14724:2024 §5.8: designation, number, dash and title ABOVE; source below.
\DeclareCaptionLabelSeparator{ft}{ – }
\captionsetup{labelsep=ft,position=top}
\DeclareFloatingEnvironment[fileext=lgr,name={Chart},listname={List of charts}]{ftchart}
\DeclareFloatingEnvironment[fileext=lqd,name={Box},listname={List of boxes}]{ftframe}
\newcommand{\ftsource}[1]{\par\smallskip{\footnotesize Source: #1\par}}
\usepackage[backend=biber,style=apa]{biblatex}
\addbibresource{bibliography.bib}
\usepackage{hyperref}
\hypersetup{hidelinks,pdftitle={Sampling and aliasing},pdfauthor={Flávio Tinoco},pdflang={en}}

% Theorem-like environments: one shared counter (amsthm \newtheorem{x}[theorem]).
\theoremstyle{plain}
\newtheorem{theorem}{Theorem}
\newtheorem{lemma}[theorem]{Lemma}
\newtheorem{proposition}[theorem]{Proposition}
\newtheorem{corollary}[theorem]{Corollary}
\newtheorem{conjecture}[theorem]{Conjecture}
\newtheorem{axiom}[theorem]{Axiom}
\theoremstyle{definition}
\newtheorem{definition}[theorem]{Definition}
\newtheorem{example}[theorem]{Example}
\theoremstyle{remark}
\newtheorem*{remark}{Remark}
\renewcommand{\qedsymbol}{\ensuremath{\blacksquare}}

% Markup of the abnt stage (semantic wrappers; see docs/guia-latex.md).
\newcommand{\ftforeign}[2]{\textit{#2}}
\newcommand{\ftname}[1]{#1}
\newcommand{\ftabbr}[2]{#1}
% KaTeX-only HTML extensions degrade to their content.
\providecommand{\htmlClass}[2]{#2}
\providecommand{\htmlData}[2]{#2}
\providecommand{\htmlId}[2]{#2}
\providecommand{\htmlStyle}[2]{#2}

% Site notation (src/core/latex/macros.ts) and document macros.
\makeatletter
\newcommand*{\ftmacro}[3]{%
  \ifdefined#1%
    \expandafter\NewCommandCopy\csname fttext\string#1\endcsname#1%
    \expandafter\newcommand\csname ftmath\string#1\endcsname[#2]{#3}%
    \DeclareRobustCommand{#1}{%
      \ifmmode\expandafter\@firstoftwo\else\expandafter\@secondoftwo\fi
      {\csname ftmath\string#1\endcsname}{\csname fttext\string#1\endcsname}}%
  \else
    \newcommand{#1}[#2]{#3}%
  \fi}
\makeatother
\ftmacro{\N}{0}{\mathbb{N}}
\ftmacro{\Z}{0}{\mathbb{Z}}
\ftmacro{\Q}{0}{\mathbb{Q}}
\ftmacro{\R}{0}{\mathbb{R}}
\ftmacro{\C}{0}{\mathbb{C}}
\ftmacro{\E}{0}{\operatorname{\mathbb{E}}}
\ftmacro{\Var}{0}{\operatorname{Var}}
\ftmacro{\Prob}{0}{\operatorname{P}}
\ftmacro{\H}{0}{\operatorname{H}}
\ftmacro{\I}{0}{\operatorname{I}}
\ftmacro{\KL}{0}{D_{\mathrm{KL}}\infdivx}
\ftmacro{\infdivx}{2}{\left(#1\,\middle\|\,#2\right)}
\ftmacro{\d}{0}{\mathop{}\!\mathrm{d}}
\ftmacro{\e}{0}{\mathrm{e}}
\ftmacro{\i}{0}{\mathrm{i}}
\ftmacro{\sinc}{0}{\operatorname{sinc}}
\ftmacro{\abs}{1}{\left\lvert #1 \right\rvert}
\ftmacro{\norm}{1}{\left\lVert #1 \right\rVert}
\ftmacro{\set}{1}{\left\{ #1 \right\}}
\ftmacro{\F}{0}{\mathcal{F}}
\ftmacro{\defeq}{0}{\coloneqq}
\ftmacro{\entails}{0}{\vDash}
\ftmacro{\proves}{0}{\vdash}

\title{Sampling and aliasing}
\author{Flávio Tinoco}
\date{October 1, 2026}

\begin{document}
\maketitle

\section{Sampling and aliasing}\label{sampling}

Sampling a signal at rate $f_s$ folds every frequency above $f_s/2$ back into the band
\parencites[{p.~11}]{shannon1949}[{see}][{}]{nyquist1928}. As \textcite[{ch.~2}]{kotelnikov1933} showed earlier, the
condition is also sufficient. Prices such as \$5 \& 10\% off, the path C:\textbackslash{}temp\textbackslash{}new and
50\_000 samples \textasciitilde{} 1 s must survive; so must x\textasciicircum{}2 written as text and a \#hashtag.

\begin{definition}[{Alias}]\label{def:alias}
A frequency $f$ is an \emph{alias} of $f + k f_s$ for every integer $k$.
\end{definition}

\begin{theorem}[{Nyquist–Shannon}]\label{thm:sampling}
A signal band-limited to $B$ is determined by its samples whenever $f_s > 2B$, see
\eqref{eq:rate} and Figure~\ref{fig:alias}.
\end{theorem}

\begin{proof}[{Proof of Theorem~\ref{thm:sampling}}]
Combine Definition~\ref{def:alias} with the reconstruction formula \eqref{eq:sinc}.
\end{proof}

\begin{equation}
f_s > 2B \label{eq:rate}
\end{equation}

\begin{align}
x(t) &= \sum_{n} x(nT)\,\sinc\!\left(\frac{t - nT}{T}\right) \label{eq:sinc}
\end{align}

\begin{gather*}
\text{an unnumbered display}
\end{gather*}

\begin{figure}[htbp]
\centering
\caption{Spectral folding}\label{fig:alias}
\includegraphics[width=\linewidth,alt={Spectrum folded around half the sampling rate}]{./alias}
\ftsource{Author’s own work (2026).}
\end{figure}

\begin{ftframe}[htbp]
\centering
\caption{Rules of thumb}\label{frm:rules}
\begin{tabular}{@{}ll@{}}
\toprule
Rule & Why \\
\midrule
Filter first & avoids \texttt{aliasing} \\
Oversample 2.5× & margin \& safety \\
\bottomrule
\end{tabular}
\ftsource{Adapted from \parencite{shannon1949}.}
\end{ftframe}

% stretch: Why a factor of two?
Because a real sinusoid occupies two points of the spectrum.

\begin{remark}
Remarks are unnumbered\footnote{like in amsthm.} and so are proofs; ratios like 3:2 stay.
\end{remark}

\begin{itemize}
\item item with [a bracket] at the start
\item \href{https://example.org/a_b?x=1\&y=2\#frag}{link} and \url{https://example.org/\~user}
\end{itemize}

\printbibliography

\end{document}
